% ==================================================================
% id:       rv-002
% topic:    משתנים מקריים
% subtopic: משתנים מקריים מוכרים
% solution: full
% verified: no
% ==================================================================
באחד מערבי חנוכה, נדב רוצה לאכול סופגנייה. בבית הוריו יש 10 סופגניות: 4 עם ריבה, 3 עם פיסטוק, 2 עם שוקולד ואחת בלי כלום. נדב אוהב רק סופגניות שוקולד, אבל כשהוא בוחר סופגנייה הוא לא רואה אותן (עוד לא הדליקו נרות, אז חשוך).

נדב בוחר באקראי סופגנייה אחת, לוקח אותה לחדר הסמוך (שם יש אור) ומסתכל עליה. אם זו לא סופגניית שוקולד, הוא זורק אותה לפח (מרוב עצבים), חוזר ובוחר באותה דרך סופגנייה נוספת מתוך הסופגניות שנשארו.

נסמן ב-$X$ את מספר הסופגניות שנדב בחר עד שקיבל סופגניית שוקולד. מהו הטווח של $X$? מצאו את פונקציית ההסתברות $P_X(k)$ לכל $k$ בטווח.
\hint{$X=k$ כשהסופגניות שנדב בוחר בפעמים $1,\dots,k-1$ אינן שוקולד, והסופגנייה שהוא בוחר בפעם ה-$k$ היא שוקולד. כפלו את ההסתברויות המותנות של כל שלב.}
\begin{solbox}
\unverified
$X$ הוא מספר הסופגניות שנדב בחר עד שקיבל סופגניית שוקולד. הערך הקטן ביותר הוא 1 (בחר שוקולד מיד), והגדול ביותר הוא 9 (בחר קודם את כל 8 הסופגניות האחרות). לכן $R_X=\set{1,\dots,9}$. נחשב ישירות:
\[
\begin{aligned}
\PP(X=1)&=\tfrac2{10},\qquad
\PP(X=2)=\tfrac8{10}\cdot\tfrac29,\qquad
\PP(X=3)=\tfrac8{10}\cdot\tfrac79\cdot\tfrac28=\tfrac{7\cdot2}{10\cdot9},\\
\PP(X=4)&=\tfrac8{10}\cdot\tfrac79\cdot\tfrac68\cdot\tfrac27=\tfrac{6\cdot2}{10\cdot9},
\end{aligned}
\]
ובאופן דומה $\PP(X=5)=\frac{5\cdot2}{90},\ \dots,\ \PP(X=8)=\frac{2\cdot2}{90}$, ולבסוף
$\PP(X=9)=\frac{8\cdot7\cdots1\cdot2}{10\cdot9\cdots2}=\frac{1\cdot2}{90}$. בנוסחה אחת: לכל $1\le k\le9$,
\[
\PP(X=k)=\frac{(10-k)\cdot2}{10\cdot9}=\ans{\frac{10-k}{45}}.
\]
\end{solbox}
